% Fragmento compartido por VII y CRISTAL; incluir despues del cuerpo Feigenbaum.
% Procedencia: resultados recuperados, con demostraciones elementales reunidas.
% Perfil, Jacobi y escalado: 00_antecedentes_operatorios.tex, 03, 06 y 07
% del presente paquete; no se cambia su selector ni su dominio uniforme.
% Constitutiva: CRISTAL_ES/main_autosuficiente.tex:58738--58840,58864--58894.
% Inversion: CRISTAL_ES/sections/cr28_restitucion_constitutiva.tex:18--59.
% Catalan: el mismo propietario, 277--373; identidad integral, unicidad y serie.
% Corte material: CUMPLIMIENTO_EDITORIAL_20260928/fuentes/CRISTAL_ES.
% No contiene una nueva calibracion ni una nueva campana de calculo.
\section{Articulation des lecteurs critique, spectral et constitutif}
\label{hmtfg:articulacion:seccion}

La composition APP--TRIT--TPK précède les lecteurs de cette section.
APP conserve dans chaque cellule les évaluations additive et multiplicative,
avec leurs restes et leurs quotients ; la réduction tritique détermine le régime,
l’orientation et la retenue de la transition. La composition effective
de la sélection, du transport et de l’actualisation de la mémoire produit l’état
enrichi et ses prolongements, développés dans
\ref{hmtfg:app-trit-transporte} et
\ref{hmtfg:subsubsec:tres-vueltas-ext-ch}.
L’holonomie nonadique compose ces transitions : le retour de phase
fait avancer la mémoire et conserve l’histoire. Les lecteurs agissent sur la
même structure discrète conjointe du continu, avec ses cinq
constructions consubstantielles et les deux lectures du même
antécédent : évaluation ordonnée et incidence de frontière.

La roue orientée et les canaux angulaires sont des sorties de cette
composition. Leurs coordonnées HMT déjà générées sont utilisées dans des opérations
ultérieures, avec leurs marques et leurs domaines ; ce ne sont pas des nombres métrologiques
introduits pour sélectionner une réponse. La lecture ordonnée
permet d’évaluer les formules suivantes, tandis que l’état de provenance
conserve la feuille, la route, le reste, le quotient, la retenue, la frontière et la mémoire.
Les restitutions démontrées retrouvent l’objet de leur domaine
déclaré ; la représentation évaluée reste accompagnée de ces
registres.

\subsection{La séparation critique comme lecture asymptotique intrinsèque}
\label{hmtfg:articulacion:separacion}

Dans le secteur uniforme de la roue dodécaphasique, les représentants
orientés et leurs poids sont
\[
 \phi_m=\frac{(3m-1)\pi_{\mathrm{HMT}}}{18},\qquad
 w_m=\frac1{12},\qquad 1\le m\le12.
\]
Le deuxième moment est normalisé par
\(\sum_{m=1}^{12}(3m-1)^2=5394=6\cdot899\).
Ainsi, le facteur angulaire commun se simplifie et il reste
\begin{equation}
 k_m=\frac{2(3m-1)}{\sqrt{899}},\qquad
 u_0(z)=\frac1{12}\sum_{m=1}^{12}\bigl(1-\cos(k_mz)\bigr),
 \qquad f_\lambda(z)=1-\lambda u_0(z).
 \label{hmtfg:articulacion:perfil}
\end{equation}
La variable \(\lambda\) parcourt la famille après fixation du
lecteur, de ses poids et de la normalisation. Dans
\(0<\lambda\le2/u_0(1)\), le domaine dynamique est \([-1,1]\).
La première racine est \(\lambda_1=1/u_0(1)\) ; pour \(n\ge2\),
\(\lambda_n\) est la plus petite racine supérieure à \(\lambda_{n-1}\)
dont le point critique a pour période exacte \(2^n\).
La classification, l’existence de chaque minimum et l’identification
métrique de ces centres ont été démontrées dans
\ref{hmtfg:escalado:15}--\ref{hmtfg:escalado:17}.

\begin{proposition}[Lecture structurelle de la constante de changement d’échelle]
\label{hmtfg:articulacion:delta}
Pour le lecteur uniforme \eqref{hmtfg:articulacion:perfil} et son sélecteur
de premières racines, la constante de séparation critique est
\begin{equation}
 \delta_{\mathrm{HMT,crit}}
 :=\lim_{n\to\infty}
 \frac{\lambda_n-\lambda_{n-1}}{\lambda_{n+1}-\lambda_n}
 =\delta_{\mathrm F}.
 \label{hmtfg:articulacion:limite}
\end{equation}
C’est l’invariant asymptotique de séparation des paramètres des retours
critiques de ce lecteur, dont l’ordre de sélection est conservé.
\end{proposition}
\begin{proof}
Le théorème de changement d’échelle \eqref{hmtfg:escalado:eq:67} fournit
\(\lambda_*-\lambda_n=C\delta_{\mathrm F}^{-n}(1+e_n)\),
avec \(C>0\) et \(e_n\to0\). En soustrayant deux termes consécutifs,
\[
 \lambda_n-\lambda_{n-1}
 =C\delta_{\mathrm F}^{-n}
   \bigl[\delta_{\mathrm F}-1+
          \delta_{\mathrm F}e_{n-1}-e_n\bigr].
\]
Le quotient de cette égalité et de celle d’indice \(n+1\) converge vers
\(\delta_{\mathrm F}\). L’identité à partir d’un certain rang des centres locaux
et des premières racines, démontrée dans
\eqref{hmtfg:escalado:eq:66}, permet d’appliquer le changement d’échelle précisément
à la suite sélectionnée, et non à une autre cascade choisie après coup.
\end{proof}

Cette formulation conserve plus d’information que le numéral terminal :
elle identifie la roue, sa lecture, la famille et le sélecteur dont la limite
est évaluée. Son interprétation est intrinsèque à cette composition HMT ;
la reconnaissance de la constante universelle utilise ensuite le théorème
analytique appliqué au lecteur déjà construit.

\subsection{La mesure finie et la restitution spectrale du profil}
\label{hmtfg:articulacion:jacobi}

Les fréquences au carré produisent la mesure de probabilité atomique
\begin{equation}
 s_m=k_m^2=\frac{4(3m-1)^2}{899},\qquad
 \nu_0=\frac1{12}\sum_{m=1}^{12}\delta_{s_m},\qquad
 M_j=\int s^j\,d\nu_0(s).
 \label{hmtfg:articulacion:medida-finita}
\end{equation}
Ici, \(\delta_{s_m}\) est la masse de Dirac en un nœud, distincte
de la constante \(\delta_{\mathrm F}\). La mesure conserve les
douze fréquences et leurs poids. Pour un polynôme non nul de degré inférieur
à \(r\le12\),
\[
 \int p(s)^2\,d\nu_0(s)
 =\frac1{12}\sum_{m=1}^{12}p(s_m)^2>0,
\]
car un polynôme de degré inférieur à douze ne peut s’annuler aux
douze nœuds distincts. Ainsi, les matrices de Hankel
\((M_{i+j})_{0\le i,j<r}\) sont définies positives.
Le polynôme \(\prod_m(s-s_m)\) est de norme nulle ; l’espace
\(L^2(\nu_0)\) est de dimension douze.

\begin{proposition}[Représentation exacte du lecteur critique]
\label{hmtfg:articulacion:representacion}
Soient \(\psi_0,\ldots,\psi_{11}\) les polynômes orthonormaux
à coefficient principal positif pour \(\nu_0\), avec
\(\psi_0=1\). La multiplication par \(s\) a pour matrice de
Jacobi \(J_{12}\), réelle symétrique à coefficients sous-diagonaux positifs.
Avec \(e_0=(1,0,\ldots,0)^{\mathsf T}\),
\begin{equation}
 u_0(z)=\int\bigl(1-\cos(z\sqrt{s})\bigr)\,d\nu_0(s)
 =e_0^{\mathsf T}\bigl[I-\cos(z\sqrt{J_{12}})\bigr]e_0.
 \label{hmtfg:articulacion:perfil-jacobi}
\end{equation}
\end{proposition}
\begin{proof}
La positivité précédente permet d’orthogonaliser jusqu’au degré onze.
Si \(i<j-1\), l’autoadjonction de la multiplication donne
\(\langle s\psi_j,\psi_i\rangle
=\langle\psi_j,s\psi_i\rangle=0\).
Par symétrie, il ne reste que la diagonale et ses deux voisines ; le rapport
positif des coefficients principaux fixe la positivité de la
sous-diagonale. Définissons
\[
 Q_{mj}=\frac{\psi_j(s_m)}{\sqrt{12}},\qquad
 D=\operatorname{diag}(s_1,\ldots,s_{12}).
\]
L’orthonormalité donne \(Q^{\mathsf T}Q=I\), et l’évaluation de la
récurrence aux nœuds donne \(DQ=QJ_{12}\). Puisque \(Q\) est
carrée, \(J_{12}=Q^{\mathsf T}DQ\). Ses valeurs propres sont
positives, de sorte que la racine positive et le cosinus se calculent dans
cette diagonalisation. Enfin,
\(Qe_0=(1,\ldots,1)^{\mathsf T}/\sqrt{12}\), ce qui prouve
\eqref{hmtfg:articulacion:perfil-jacobi}.
\end{proof}

La réalisation par \((J_{12},e_0)\) restitue le même profil ;
ainsi, en conservant la famille et le sélecteur, elle restitue également
ses retours et ses premières racines. La positivité spectrale
justifie cette représentation. La transversalité paramétrique
utilise en outre la dérivée du retour et ses produits orbitaux,
développés dans \ref{hmtfg:escalado:14.3} et
\ref{hmtfg:escalado:6}. Ce sont des propriétés d’opérations différentes :
la démonstration du changement d’échelle conserve les deux et ne remplace pas le
contrôle de la tangente par la positivité des moments.

\subsection{Les canaux constitutifs et leur inversion étiquetée}
\label{hmtfg:articulacion:constitutiva}

Le lecteur constitutif agit sur un autre objet typé du même
support : le transport angulaire à deux feuilles. Ses coordonnées déjà
générées sont \(x=A_{\rm rad}\) et \(y=C^*_{\rm rad}\), dans la
chambre contractante \(x>|y|\). L’involution des feuilles
\(R=R^*=R^{-1}\) détermine les projecteurs marqués
\(P_\pm=(I\pm R)/2\). La représentation du transport est
\begin{equation}
 q_+=\exp[-(x+y)],\qquad q_-=\exp[-(x-y)],\qquad
 T=q_+P_++q_-P_-,\qquad 0<q_\pm<1.
 \label{hmtfg:articulacion:canales}
\end{equation}
L’évaluation exponentielle réalise ces canaux après leur
construction angulaire. Les incidences \(90\) et \(120\) du
transport fixent la réponse
\[
 \mathcal V=(I-T^{90})(I-T^{120})^{-1}.
\]
L’inverse existe parce que \(1-q_\pm^{120}>0\).
En écrivant \(s=q^{30}\), avec \(30=\gcd(90,120)\), on obtient
\begin{equation}
 f(s)=\frac{1-s^3}{1-s^4}
     =\frac{1+s+s^2}{1+s+s^2+s^3},\qquad
 r_\pm=f(q_\pm^{30}),\qquad
 \mathcal V=r_+P_++r_-P_-.
 \label{hmtfg:articulacion:respuesta}
\end{equation}
Dans cette sous-section, \(f\) désigne la réponse constitutive ;
\(f_\lambda\) continue de désigner l’application dynamique de
\eqref{hmtfg:articulacion:perfil}. Leurs arguments et leurs domaines
sont distincts.

\begin{proposition}[Restitution des canaux et de l’orientation]
\label{hmtfg:articulacion:inversion}
La fonction constitutive \(f:(0,1)\to(3/4,1)\) est une bijection
décroissante. Son inverse \(\psi\) associe à \(r\) l’unique racine
\(s\in(0,1)\) de
\begin{equation}
 r s^3+(r-1)(1+s+s^2)=0.
 \label{hmtfg:articulacion:cubica}
\end{equation}
Les lectures constitutives normalisées, avec leurs feuilles étiquetées,
\[
 \widehat\varepsilon=r_+^2,\qquad
 \widehat\mu=r_-^2,\qquad
 \widehat Z=\frac{r_-}{r_+},\qquad
 \widehat c=\frac1{r_+r_-},
\]
avec \(\widehat\varepsilon,\widehat\mu\in(9/16,1)\), permettent de restituer
\begin{equation}
 \begin{aligned}
 r_+&=\sqrt{\widehat\varepsilon},&
 r_-&=\sqrt{\widehat\mu},\\
 s_\pm&=\psi(r_\pm),& q_\pm&=s_\pm^{1/30}.
 \end{aligned}
 \label{hmtfg:articulacion:recuperacion-canales}
\end{equation}
\begin{equation}
 x=-\frac12\log(q_+q_-),\qquad
 y=\frac12\log\frac{q_-}{q_+}.
 \label{hmtfg:articulacion:recuperacion-angular}
\end{equation}
\end{proposition}
\begin{proof}
La dérivation du quotient rationnel donne
\begin{equation}
 f'(s)=-\frac{s^2(3+2s+s^2)}{(1+s+s^2+s^3)^2}<0
 \qquad(0<s<1).
 \label{hmtfg:articulacion:monotonia}
\end{equation}
Les limites \(f(0^+)=1\) et \(f(1^-)=3/4\), jointes à la
continuité, prouvent la bijectivité. Multiplier \(f(s)=r\) par
le dénominateur positif produit \eqref{hmtfg:articulacion:cubica}.
La positivité de \(r_\pm\) détermine les racines carrées
et celle de \(s_\pm\) détermine les racines trentièmes.
Par \eqref{hmtfg:articulacion:canales},
\(\log q_+=-x-y\) et \(\log q_-=-x+y\) ;
l’addition et la soustraction restituent \(x,y\).
Réciproquement, les racines restituées appartiennent à \((0,1)\),
de sorte que \(x\pm y=-\log q_\pm>0\) ; la paire obtenue
reste dans la chambre \(x>|y|\).
\end{proof}

En particulier,
\(\widehat\mu\widehat\varepsilon=\widehat c^{-2}\) et
\(\widehat\mu/\widehat\varepsilon=\widehat Z^2\).
L’échange des étiquettes permute \(q_+\) et \(q_-\),
conserve \(x\) et change le signe de \(y\). La restitution
exige cette orientation : le produit isolé \(r_+r_-\)
conserve la propagation commune, mais ne détermine pas les deux réponses
étiquetées. Les racines positives et l’inversion cubique restituent
les canaux à l’intérieur du lecteur fixé ; le registre complet de leur
génération reste dans l’état enrichi.

\subsection{Le moment de Catalan et la restitution fonctionnelle}
\label{hmtfg:articulacion:catalan}

Les mêmes incidences déterminent une relation entre fonctions,
avant l’évaluation de l’un ou l’autre canal. Pour \(0<s<1\), soit
\(\mathcal D=s\,d/ds\), et définissons le moment orienté
\begin{equation}
 \mathcal C_2(s)
 =\int_0^s\log\frac{s}{t}\,
     \frac{1+t}{1+t+t^2}f(t)\,dt
 =\int_0^s\frac{\log(s/t)}{1+t^2}\,dt.
 \label{hmtfg:articulacion:integral}
\end{equation}
La deuxième égalité découle de la factorisation
\(1+t+t^2+t^3=(1+t)(1+t^2)\). Le poids continu et positif
\(dt/(1+t^2)\) est ainsi produit par la composition rationnelle
du lecteur constitutif, et non par une identification à
\(\nu_0\).

\begin{proposition}[Inversion différentielle, unicité et lecture à l’extrémité]
\label{hmtfg:articulacion:restitucion-catalan}
Le moment \(\mathcal C_2\) est l’unique fonction
\(F\in C^2((0,1))\) qui satisfait
\begin{equation}
 \mathcal D^2F(s)=\frac{s(1+s)}{1+s+s^2}f(s),\qquad
 \lim_{s\downarrow0}F(s)=0.
 \label{hmtfg:articulacion:problema-diferencial}
\end{equation}
La réponse constitutive est restituée par
\begin{equation}
 \mathcal D^2\mathcal C_2(s)=\frac{s}{1+s^2},\qquad
 f(s)=\frac{1+s+s^2}{s(1+s)}\,
          \mathcal D^2\mathcal C_2(s).
 \label{hmtfg:articulacion:inversa-diferencial}
\end{equation}
Son extension continue à l’intervalle fermé possède la série
\begin{equation}
 \mathcal C_2(s)=\sum_{j=0}^{\infty}
       \frac{(-1)^j s^{2j+1}}{(2j+1)^2},\qquad 0\le s\le1,
 \qquad
 \mathcal C_2(1)=\sum_{j=0}^{\infty}\frac{(-1)^j}{(2j+1)^2}
 =G_{\rm Cat}.
 \label{hmtfg:articulacion:serie-catalan}
\end{equation}
\end{proposition}
\begin{proof}
L’intégrande de \eqref{hmtfg:articulacion:integral} est positif ou nul
et
\[
 0\le\mathcal C_2(s)\le\int_0^s\log(s/t)\,dt=s.
\]
Cela prouve l’intégrabilité en zéro et la limite requise.
L’identité \(\log(s/t)=\int_t^s du/u\), suivie de
l’interversion d’intégrales positives, fournit
\[
 \mathcal C_2(s)
 =\int_0^s\frac1u\left(\int_0^u\frac{dt}{1+t^2}\right)du.
\]
Sur tout intervalle compact du domaine ouvert, on applique le
théorème fondamental du calcul intégral. Par conséquent,
\[
 \mathcal D\mathcal C_2(s)=\int_0^s\frac{dt}{1+t^2},
 \qquad
 \mathcal D^2\mathcal C_2(s)=\frac{s}{1+s^2}.
\]
La factorisation rationnelle utilisée dans
\eqref{hmtfg:articulacion:integral} donne
\eqref{hmtfg:articulacion:problema-diferencial} ; isoler \(f\)
donne \eqref{hmtfg:articulacion:inversa-diferencial}, dont les
dénominateurs sont positifs dans \((0,1)\).

Si deux solutions ont la limite nulle requise, leur différence
\(H\) satisfait \(\mathcal D^2H=0\). Le changement \(z=\log s\)
transforme cette équation en \(d^2H(e^z)/dz^2=0\), d’où
\(H(s)=a\log s+b\). L’existence de la limite nulle impose
\(a=b=0\) et prouve l’unicité.

Pour \(s<1\), le développement géométrique de \((1+t^2)^{-1}\)
peut être intégré terme à terme avec le poids logarithmique, car
\[
 \int_0^s t^{2j}\log(s/t)\,dt
 =\frac{s^{2j+1}}{(2j+1)^2},\qquad
 \sum_{j\ge0}\frac{s^{2j+1}}{(2j+1)^2}<\infty.
\]
Cela prouve la série à l’intérieur. La borne
\(1/(2j+1)^2\) donne la convergence absolue et uniforme de cette
série sur \([0,1]\). Dans l’intégrale, l’intégrande prolongé
par zéro pour \(t>s\) est dominé par \(\log(1/t)\)
sur \((0,1)\). La convergence dominée permet d’atteindre
\(s=1\) et coïncide avec le prolongement de la série. Sa valeur
à l’extrémité est précisément la série de Catalan indiquée.
\end{proof}

L’inversion différentielle utilise la famille \(\mathcal C_2(s)\),
et non seulement sa valeur à l’extrémité \(G_{\rm Cat}\). En particulier,
\(\mathcal D^2\mathcal C_2(s)\to1/2\) lorsque \(s\uparrow1\).
La série dérivée deux fois avec \(\mathcal D\) y possède
une limite radiale d’Abel ; la série du moment lui-même est
absolument convergente à la frontière. Cette distinction maintient
les conditions précises de chaque passage à la limite.

\subsection{Composition des restitutions et portée conjointe}
\label{hmtfg:articulacion:conclusion}

Les interrelations précédentes ont des opérations et des domaines
explicites. Les moments de la mesure atomique \(\nu_0\)
produisent \(J_{12}\), et son calcul fonctionnel restitue le profil
\(u_0\). Le profil et le sélecteur de premières racines déterminent
la suite dont le changement d’échelle fournit la lecture \(\delta_{\mathrm F}\).
D’autre part, la fonction constitutive \(f\) produit
\(\mathcal C_2\) au moyen du moment intégral, et
\eqref{hmtfg:articulacion:inversa-diferencial} restitue \(f\)
à partir de cette famille. Ses évaluations étiquetées en
\(s_\pm=q_\pm^{30}\) déterminent les réponses constitutives ;
l’inversion cubique restitue les deux canaux et leurs coordonnées
angulaires. La lecture du moment à l’extrémité donne \(G_{\rm Cat}\).

La mesure \(\nu_0\) est positive, atomique et de rang douze.
Le moment de Catalan intègre un poids positif continu et possède
un développement alterné. Les signes de ce développement sont des
coefficients de la série géométrique ; ce ne sont pas des poids négatifs de
\(\nu_0\), et ils ne convertissent pas l’intégrale positive en la même
représentation spectrale. Chaque lecteur conserve son support,
son orientation, sa normalisation et son opération d’évaluation.

L’unité réside dans la structure APP--TRIT--TPK qui engendre et
transporte leurs antécédents ; les relations fonctionnelles sont
établies par les compositions démontrées. L’égalité
d’origine n’identifie pas des opérateurs de domaines différents.
En particulier, cette articulation ne postule pas de formule scalaire
qui détermine \(\delta_{\mathrm F}\) à partir de la constante de Catalan et
d’autres constantes physiques. Elle situe le changement d’échelle critique aux côtés de
restitutions effectives entre lecteurs du même support,
en conservant l’information requise par chaque direction.
